% GENERATED FILE. Edit canonical lessons, then run npm run generate:books.
% canonical-source-hash: fnv1a64:a089532a772f0238
% canonical-lessons: TE-C483-aviravu, TE-C483-inkipovu, TE-C483-mandu, TE-C483-aripovu, TE-C483-cigurincu

\chapter{Water, Fire and Light}
\label{ch:te-water-fire-and-light}
\hlchaptermodality{483}

\begin{chapteropening}
\textbf{By the end of this chapter:} \emph{I can say to evaporate, to soak in, to blaze, to go out and to put forth leaves.}

Complete the last lesson of chapter 483: I can say to evaporate, to soak in, to blaze, to go out and to put forth leaves.
\end{chapteropening}

\section[\texorpdfstring{\te{ఆవిరవు} (āviravu)}{ఆవిరవు (āviravu)}]{\te{ఆవిరవు} (āviravu) — to evaporate}

\label{lesson:TE-C483-aviravu}



\begin{quote}
Before the new one: say the Telugu for a breeze, then the Telugu for to cool down.
\end{quote}

\subsection*{You'll want to know: \te{ఆవిరవు}}
\textbf{\te{ఆవిరవు}} — \emph{āviravu} — \textquotedblleft{}to evaporate\textquotedblright{}.

\begin{quote}
\te{ఎండకు} \te{నేల} \te{మీద} \te{నీరు} \te{ఆవిరైంది}.
\end{quote}

\begin{quote}
\emph{eṇḍaku nēla mīda nīru āviraindi.}
\end{quote}

\begin{quote}
In the sun the water on the ground evaporated.
\end{quote}

Past (it …): \te{ఆవిరయింది}, \emph{āvirayindi}.

\te{ఆవిరి}, steam, + \te{అవు}, to become. A hope that vanishes also \te{ఆవిరవుతుంది}. Its subject is a thing, so the past shown is the third person: \te{ఆవిరయింది}, it evaporated.

\begin{cousinweb}
\textbf{\te{ఆవిరవు}} (\emph{āviravu}) is native Dravidian: \te{ఆవిరి}, \textquotedblleft{}vapour, steam\textquotedblright{} (DEDR 393), with \te{అవు}, \textquotedblleft{}to become\textquotedblright{}.
\end{cousinweb}

\begin{grammarlens}[title={What you've built}]
One more word for getting around.
\end{grammarlens}

\subsection*{Your turn}
\begin{itemize}
\raggedright
  \item \emph{Say it:} \emph{āviravu}
  \item \emph{Say it:} \emph{āviravu}, once more
  \item \emph{Say it:} \emph{eṇḍaku nēla mīda nīru āviraindi.} — In the sun the water on the ground evaporated.
  \item \emph{Recall:} say the Telugu for a breeze, then the Telugu for to cool down, then say \emph{āviravu} again
\end{itemize}

\begin{tcolorbox}[breakable,colback=teal!4,colframe=teal!35!black,title={Before you move on}]
What does \te{ఆవిరవు} mean? (\textquotedblleft{}To evaporate\textquotedblright{}.) Say it once more.
\end{tcolorbox}
\section[\texorpdfstring{\te{ఇంకిపోవు} (iṅkipōvu)}{ఇంకిపోవు (iṅkipōvu)}]{\te{ఇంకిపోవు} (iṅkipōvu) — to soak in; to sink in}

\label{lesson:TE-C483-inkipovu}



\begin{quote}
Before the new one: say the Telugu for to cool down, then the Telugu for to evaporate.
\end{quote}

\subsection*{You'll want to know: \te{ఇంకిపోవు}}
\textbf{\te{ఇంకిపోవు}} — \emph{iṅkipōvu} — \textquotedblleft{}to soak in; to sink in\textquotedblright{}.

\begin{quote}
\te{వర్షం} \te{నీరు} \te{నేలలో} \te{ఇంకిపోయింది}.
\end{quote}

\begin{quote}
\emph{varṣaṁ nīru nēlalō iṅkipōyindi.}
\end{quote}

\begin{quote}
The rainwater soaked into the ground.
\end{quote}

Past (it …): \te{ఇంకిపోయింది}, \emph{iṅkipōyindi}.

Said of water going into the earth, and of a pond drying up. Its subject is a thing, so the past shown is the third person: \te{ఇంకిపోయింది}, it soaked in.

\begin{cousinweb}
\textbf{\te{ఇంకిపోవు}} (\emph{iṅkipōvu}) is native Dravidian: DEDR 430 lists Telugu \emph{iṅku}, \textquotedblleft{}to dry up, to sink in\textquotedblright{}, here with \te{పోవు}, \textquotedblleft{}to go\textquotedblright{} (DEDR 4572).
\end{cousinweb}

\begin{grammarlens}[title={What you've built}]
One more word for getting around.
\end{grammarlens}

\subsection*{Your turn}
\begin{itemize}
\raggedright
  \item \emph{Say it:} \emph{iṅkipōvu}
  \item \emph{Say it:} \emph{iṅkipōvu}, once more
  \item \emph{Say it:} \emph{varṣaṁ nīru nēlalō iṅkipōyindi.} — The rainwater soaked into the ground.
  \item \emph{Recall:} say the Telugu for to cool down, then the Telugu for to evaporate, then say \emph{iṅkipōvu} again
\end{itemize}

\begin{tcolorbox}[breakable,colback=teal!4,colframe=teal!35!black,title={Before you move on}]
What does \te{ఇంకిపోవు} mean? (\textquotedblleft{}To soak in; to sink in\textquotedblright{}.) Say it once more.
\end{tcolorbox}
\section[\texorpdfstring{\te{మండు} (maṇḍu)}{మండు (maṇḍu)}]{\te{మండు} (maṇḍu) — to blaze; to burn}

\label{lesson:TE-C483-mandu}



\begin{quote}
Before the new one: say the Telugu for to evaporate, then the Telugu for to soak in.
\end{quote}

\subsection*{You'll want to know: \te{మండు}}
\textbf{\te{మండు}} — \emph{maṇḍu} — \textquotedblleft{}to blaze; to burn\textquotedblright{}.

\begin{quote}
\te{మధ్యాహ్నం} \te{ఎండ} \te{మండుతోంది}.
\end{quote}

\begin{quote}
\emph{madhyāhnaṁ eṇḍa maṇḍutōndi.}
\end{quote}

\begin{quote}
At midday the sun is blazing.
\end{quote}

Past (it …): \te{మండింది}, \emph{maṇḍindi}.

Said of fire, of strong sun and of a burning feeling: \te{కళ్ళు} \te{మండుతున్నాయి}, my eyes are burning. Its subject is a thing, so the past shown is the third person: \te{మండింది}, it blazed.

\begin{cousinweb}
\textbf{\te{మండు}} (\emph{maṇḍu}) is native Dravidian: DEDR 4680 lists Telugu \emph{maṇḍu}, \textquotedblleft{}to burn, to blaze, to flame\textquotedblright{}, with Tamil \emph{maṇṭu}, \textquotedblleft{}to blaze up\textquotedblright{}.
\end{cousinweb}

\begin{grammarlens}[title={What you've built}]
One more word for getting around.
\end{grammarlens}

\subsection*{Your turn}
\begin{itemize}
\raggedright
  \item \emph{Say it:} \emph{maṇḍu}
  \item \emph{Say it:} \emph{maṇḍu}, once more
  \item \emph{Say it:} \emph{madhyāhnaṁ eṇḍa maṇḍutōndi.} — At midday the sun is blazing.
  \item \emph{Recall:} say the Telugu for to evaporate, then the Telugu for to soak in, then say \emph{maṇḍu} again
\end{itemize}

\begin{tcolorbox}[breakable,colback=teal!4,colframe=teal!35!black,title={Before you move on}]
What does \te{మండు} mean? (\textquotedblleft{}To blaze; to burn\textquotedblright{}.) Say it once more.
\end{tcolorbox}
\section[\texorpdfstring{\te{ఆరిపోవు} (āripōvu)}{ఆరిపోవు (āripōvu)}]{\te{ఆరిపోవు} (āripōvu) — to go out; to die down}

\label{lesson:TE-C483-aripovu}



\begin{quote}
Before the new one: say the Telugu for to soak in, then the Telugu for to blaze.
\end{quote}

\subsection*{You'll want to know: \te{ఆరిపోవు}}
\textbf{\te{ఆరిపోవు}} — \emph{āripōvu} — \textquotedblleft{}to go out; to die down\textquotedblright{}.

\begin{quote}
\te{గాలికి} \te{దీపం} \te{ఆరిపోయింది}.
\end{quote}

\begin{quote}
\emph{gāliki dīpaṁ āripōyindi.}
\end{quote}

\begin{quote}
The lamp went out in the wind.
\end{quote}

Past (it …): \te{ఆరిపోయింది}, \emph{āripōyindi}.

\te{ఆరు}, to cool, to go out, + \te{పోవు}. A fire, a lamp or wet clothes drying: \te{బట్టలు} \te{ఆరిపోయాయి}. Its subject is a thing, so the past shown is the third person: \te{ఆరిపోయింది}, it went out.

\begin{cousinweb}
\textbf{\te{ఆరిపోవు}} (\emph{āripōvu}) is native Dravidian: \te{ఆరు}, \textquotedblleft{}to cool, to be put out, to dry\textquotedblright{} (DEDR 404), with \te{పోవు}, \textquotedblleft{}to go\textquotedblright{} (DEDR 4572).
\end{cousinweb}

\begin{grammarlens}[title={What you've built}]
One more word for getting around.
\end{grammarlens}

\subsection*{Your turn}
\begin{itemize}
\raggedright
  \item \emph{Say it:} \emph{āripōvu}
  \item \emph{Say it:} \emph{āripōvu}, once more
  \item \emph{Say it:} \emph{gāliki dīpaṁ āripōyindi.} — The lamp went out in the wind.
  \item \emph{Recall:} say the Telugu for to soak in, then the Telugu for to blaze, then say \emph{āripōvu} again
\end{itemize}

\begin{tcolorbox}[breakable,colback=teal!4,colframe=teal!35!black,title={Before you move on}]
What does \te{ఆరిపోవు} mean? (\textquotedblleft{}To go out; to die down\textquotedblright{}.) Say it once more.
\end{tcolorbox}
\section[\texorpdfstring{\te{చిగురించు} (ciguriñcu)}{చిగురించు (ciguriñcu)}]{\te{చిగురించు} (ciguriñcu) — to put forth leaves; to sprout}

\label{lesson:TE-C483-cigurincu}



\begin{quote}
Before the new one: say the Telugu for to blaze, then the Telugu for to go out.
\end{quote}

\subsection*{You'll want to know: \te{చిగురించు}}
\textbf{\te{చిగురించు}} — \emph{ciguriñcu} — \textquotedblleft{}to put forth leaves; to sprout\textquotedblright{}.

\begin{quote}
\te{వర్షాలు} \te{పడగానే} \te{చెట్లు} \te{చిగురించాయి}.
\end{quote}

\begin{quote}
\emph{varṣālu paḍagānē ceṭlu ciguriñcāyi.}
\end{quote}

\begin{quote}
As soon as the rains came the trees put out new leaves.
\end{quote}

Past (it …): \te{చిగురించింది}, \emph{ciguriñcindi}.

\te{చిగురు}, a tender shoot, + -\te{ఇంచు}. Hope also \te{చిగురిస్తుంది}. Its subject is a thing, so the past shown is the third person: \te{చిగురించింది}, it put out leaves.

\begin{cousinweb}
\textbf{\te{చిగురించు}} (\emph{ciguriñcu}) is native Dravidian: DEDR 2489 lists Telugu \emph{iguru}, \textquotedblleft{}a bud, a sprout\textquotedblright{}, and \emph{igirincu}, \textquotedblleft{}to bud\textquotedblright{}; \te{చిగురు} is its everyday form.
\end{cousinweb}

\begin{grammarlens}[title={What you've built}]
One more word for getting around.
\end{grammarlens}

\subsection*{Your turn}
\begin{itemize}
\raggedright
  \item \emph{Say it:} \emph{ciguriñcu}
  \item \emph{Say it:} \emph{ciguriñcu}, once more
  \item \emph{Say it:} \emph{varṣālu paḍagānē ceṭlu ciguriñcāyi.} — As soon as the rains came the trees put out new leaves.
  \item \emph{Recall:} say the Telugu for to blaze, then the Telugu for to go out, then say \emph{ciguriñcu} again
\end{itemize}

\begin{tcolorbox}[breakable,colback=teal!4,colframe=teal!35!black,title={Before you move on}]
What does \te{చిగురించు} mean? (\textquotedblleft{}To put forth leaves; to sprout\textquotedblright{}.) Say it once more.
\end{tcolorbox}
